\section{方法}
\label{sec:method}

\subsection{持续学习复合目标}

在标准GRPO策略梯度之外叠加三项持续学习正则，构成复合损失：

\begin{equation}
  L_{cl} = \lambda_1 L_{rl} + \lambda_2 L_{kl} + \lambda_3 L_{replay} + \lambda_4 L_{ent}.
  \label{eq:cl-loss}
\end{equation}
$L_{rl}$为GRPO项（组内奖励归一化估计优势）。$L_{kl}=D_{KL}(\pi_{new}\,\|\,\pi_{ref})$为对参考策略$\pi_{ref}$的逆向KL锚定。$L_{replay}$为经验回放损失，作用于从缓冲区采样的历史轨迹（第\ref{sec:buffer}节）；回放行携带独立的\texttt{replay\_response\_mask}，不参与PPO主损失。$L_{ent}$为熵正则，$\lambda_4=0.001$所有实验固定开启，防策略多样性坍缩~\cite{B4}。早期参数L2正则$L_{reg}$已弃用。目标以\textbf{零框架改动}注入——通过框架的损失注入接口与钩子系统挂载。

\subsection{奖励设计}
\label{sec:reward}

由单个外部冻结大模型裁判在统一细则下打分，四个维度：

\begin{itemize}
  \item \textbf{Done}（完成度，$0/1$）：任务是否真正完成。
  \item \textbf{Correctness}（正确性，$[0,1]$）：产出是否正确。
  \item \textbf{Trajectory}（轨迹质量，$[0,1]$）：工具调用质量、有无冗余步、推理是否连贯；\textbf{并惩罚"加了不该加的东西"}。
  \item \textbf{Safety}（安全，\textbf{二元$0/1$}）：是否出现破坏性/越权操作。二元门控而非连续打分，一致性更高。
\end{itemize}

最终奖励乘性聚合：

\begin{equation}
  r =
  \begin{cases}
    (0.4\cdot\text{correctness} + 0.4\cdot\text{trajectory} + 0.2)\cdot\text{safety}, & \text{done}=1,\\[4pt]
    0.4\cdot\text{trajectory}\cdot\text{safety}, & \text{done}=0.
  \end{cases}
  \label{eq:reward}
\end{equation}

安全作为$\{0,1\}$门控，一旦触发整条清零。未完成任务仍给轨迹分，但无完成基线分和正确性分。

\subsection{九桶经验回放缓冲区}
\label{sec:buffer}

\paragraph{按能力分桶。}桶按能力/领域划分。九个纯文本能力域（表\ref{tab:buckets}）由评测基准官方类目合并、去多模态得到。$K=9$桶，共195个纯文本任务，单层无子桶。

\begin{table}[t]
\centering
\caption{各桶任务数与容量上限。}
\label{tab:buckets}
\begin{tabular}{lrr}
\toprule
桶（能力域） & 任务数 & 容量上限 \\
\midrule
workflow（工作流编排） & 56 & 4915 \\
ops（系统操作） & 44 & 4354 \\
qa（问答检索） & 36 & 3938 \\
finance（财务金融） & 20 & 2935 \\
office（办公文档） & 11 & 2177 \\
communication（沟通表达） & 11 & 2177 \\
safety（安全合规） & 9 & 1969 \\
research（研究综合） & 6 & 1607 \\
coding（代码） & 2 & 928 \\
\midrule
合计 & 195 & 25000 \\
\bottomrule
\end{tabular}
\end{table}

\paragraph{容量分配。}各桶容量$\text{cap}_i$由任务数$n_i$的平方根加权分配，总容量$C=25000$：

\begin{equation}
  \text{cap}_i = C\cdot\frac{n_i^{\alpha}}{\sum_j n_j^{\alpha}},\qquad \alpha=0.5.
  \label{eq:quota}
\end{equation}
取$\alpha=0.5$（次线性）使大桶温和倾斜、小桶被抬起。每桶另设保护下限（约为cap的$1/30$）。

\paragraph{淘汰与采样。}淘汰\textbf{仅在桶内进行}——禁止跨桶挤出。默认淘汰策略为FIFO。CLEAR基线用蓄水池采样~\cite{A2}。回放采样\textbf{两级}：先按能力空间\textbf{距离}选桶，离当前训练批质心越远权重越大（遗忘风险越高）；桶内随机（uniform）采样。每步回放行总数固定。

\paragraph{冷启动种子。}预采集的完整轨迹（含工具调用与思考）可预填缓冲区。

\subsection{回放损失加权}
\label{sec:replay-weight}

每条回放轨迹在$L_{replay}$中按token级权重：

\begin{equation}
  w_t^{(i)} = \text{normalize}\!\Big(\text{clip}\big(\text{priority}_i \cdot \tfrac{\gamma^{\text{block}(t)} + \delta^{K_i - \text{block}(t)}}{2},\; q_5,\; q_{95}\big)\Big),
  \label{eq:u-weight}
\end{equation}
其中块按OpenAI chat message边界划分。第二个因子为U形块权重的一般形式。

\textbf{当前默认关闭U形：}$\gamma=\delta=1$，块权重退化为常数1.0，退化为纯priority加权（scheme W0）。U形（$\gamma=\delta<1$，scheme W2）作为独立消融项保留。
